% !PRINTED none
% Abschnitt (part) title page: "Vierzehnter Abschnitt." / section title, rule, then opening of §.152 with its heading and the start of the section's running text and first footnote. No running-head page number is printed on this title page.
\begin{center}
\gesperrt{Vierzehnter Abschnitt.}
\end{center}

\begin{center}
\gesperrt{Anwendung der Permutationsgruppen auf Gleichungen.}
\end{center}

\begin{center}
\rule{2cm}{0.4pt}
\end{center}

\begin{center}
\S.\,152.
\end{center}

\begin{center}
\gesperrt{Wirkung der Permutationsgruppen auf Functionen von unabhängigen Veränderlichen.}
\end{center}

Für ein tieferes Eindringen in die Algebra ist nach den Ergebnissen des vorigen Abschnittes ein genaueres Studium der Permutationsgruppen erforderlich. Wir stellen zunächst einige allgemeine Sätze darüber auf, die für das Folgende die Grundlage bilden\footnote{Für die Theorie der Permutationen sind besonders hervorzuheben: \gesperrt{Cauchy}, „Journal de l'École polytechn. X. cah. (1815)“ (mehrere Abhandlungen). „C. \gesperrt{Jordan}, Traité des substitutions et des équations algébriques“, Paris 1870. \gesperrt{Netto}, „Substitutionentheorie und ihre Anwendung auf die Algebra.“ Leipzig 1882.}.

Es sei
\weq{1}{u,\ u_1,\ u_2 \ldots u_{m-1}}
ein System von einander unabhängiger Zeichen (Veränderliche) und
\weq{2}{\psi(u,\ u_1,\ u_2 \ldots u_{m-1})}
eine ganze rationale Function von ihnen mit Coëfficienten aus einem beliebigen Körper.

Als gleich sind zwei solche Functionen $\psi$ nur dann zu betrachten, wenn in den nach Potenzen und Producten der $u$ geordneten Ausdrücken entsprechende Glieder die gleichen Coëfficienten haben.
